%!TEX root = forallxyyc.tex

\chapter{Sistemas de prova alternativos}
Ao formular nosso sistema de dedução natural, tratamos certas regras de dedução natural como \emph{básicas} e outras como \emph{derivadas}. No entanto, poderíamos igualmente ter tomado várias regras diferentes como básicas ou derivadas. Ilustraremos esse ponto considerando tratamentos alternativos da disjunção, da negação e dos quantificadores. Também explicaremos por que fizemos nossas escolhas.


\section{Eliminação alternativa da disjunção}
Alguns sistemas tomam DS como sua regra básica para a eliminação da disjunção. Esses sistemas podem então tratar a regra $\eor$E como uma regra derivada. Para isso, poderiam oferecer o seguinte esquema de prova:
\begin{fitchproof}
  \have[m]{ab}{\metav{A}\eor\metav{B}}
  \open
    \hypo[i]{a}{\metav{A}} \AS
    \have[j]{c1}{\metav{C}}
  \close
  \open
    \hypo[k]{b}{\metav{B}} \AS
    \have[l]{c2}{\metav{C}}
  \close
  \have[n]{aic}{\metav{A} \eif \metav{C}}\ci{a-c1}
  \have{bic}{\metav{B} \eif \metav{C}}\ci{b-c2}
  \open
    \hypo{nc}{\enot\metav{C}} \AS
    \open
      \hypo{a2}{\metav{A}} \AS
      \have{c3}{\metav{C}}\ce{a2,aic}
      \have{bot1}{\ered}\ne{nc,c3}
    \close
    \have{na}{\enot\metav{A}}\ni{a2-bot1}
    \have{b2}{\metav{B}}\ds{ab,na}
    \have{c4}{\metav{C}}\ce{b2,bic}
    \have{bot2}{\ered}\ne{nc,c4}
  \close
  \have{con}{\metav{C}}\ip{nc-bot2}
\end{fitchproof}
Então, por que escolhemos tomar $\eor$E como básica, em vez de DS?\footnote{A versão original deste livro, de P.~D.\ Magnus, fazia o contrário.} Nosso raciocínio é que DS envolve o uso de \textquotedblleft $\enot$\textquotedblright{} no enunciado da regra. Em certo sentido, é \textquotedblleft mais limpa\textquotedblright{} nossa regra de eliminação da disjunção quando ela evita mencionar \emph{outros} conectivos.


\section{Regras alternativas para a negação}
Alguns sistemas tomam a seguinte regra como sua regra básica de introdução da negação:
\begin{fitchproof}
	\open
		\hypo[m]{a}{\metav{A}}\AS
		\have[n-1]{b}{\metav{B}}
		\have[n]{nb}{\enot\metav{B}}
	\close
	\have[\ ]{}{\enot\metav{A}}\by{$\enot$I*}{a-nb}
\end{fitchproof}
e uma versão correspondente da regra que chamamos de IP como sua regra básica de eliminação da negação:
\begin{fitchproof}
	\open
		\hypo[m]{na}{\enot\metav{A}} \AS
		\have[n][-1]{b}{\metav{B}}
		\have{nb}{\enot\metav{B}}
	\close
	\have[\ ]{a}{\metav{A}}\by{$\enot$E*}{na-nb}
\end{fitchproof}
Usando essas duas regras, poderíamos ter evitado completamente o uso do símbolo \textquotedblleft $\ered$\textquotedblright{}.\footnote{Novamente, a versão original deste livro, de P.~D.\ Magnus, fazia o contrário.} O sistema resultante teria menos regras que o nosso.

Outra maneira de lidar com a negação é usar LEM ou DNE como regra básica e introduzir IP como regra derivada. Normalmente, em um sistema desse tipo, as regras também recebem nomes diferentes. Por exemplo, às vezes o que chamamos de $\enot$E é chamado de $\ered$I, e o que chamamos de X é chamado de $\ered$E.\footnote{A versão deste livro devida a Tim Button segue esse caminho e substitui IP por LEM, que ele chama de TND, de \textit{tertium non datur}.}

Então, por que escolhemos nossas regras para a negação e a contradição?

Nossa primeira razão é que acrescentar o símbolo \textquotedblleft $\ered$\textquotedblright{} ao nosso sistema de dedução natural torna as provas consideravelmente mais fáceis de manipular. Por exemplo, em nosso sistema, é sempre claro qual é a conclusão de uma subprova: a sentença na última linha, por exemplo, $\ered$ em IP ou $\enot$I. Em $\enot$I* e $\enot$E*, as subprovas têm duas conclusões, de modo que não é possível verificar de relance se uma aplicação delas está correta.

Nossa segunda razão é que uma grande quantidade de discussão filosófica fascinante se concentrou na aceitabilidade, ou não, da prova indireta IP (equivalentemente, do terceiro excluído, isto é, LEM, ou da eliminação da dupla negação, DNE) e da explosão (isto é, X). Ao tratá-las como regras separadas no sistema de provas, você estará em melhor posição para participar dessa discussão filosófica. Em particular: depois de invocar essas regras explicitamente, será muito mais fácil saber como seria um sistema que não tivesse essas regras.

Essa discussão — e, na verdade, a grande maioria dos estudos matemáticos sobre aplicações das provas de dedução natural além dos cursos introdutórios — faz referência a uma versão diferente da dedução natural. Essa versão foi inventada por Gerhard Gentzen em 1935 e refinada por Dag Prawitz em 1965. Nosso conjunto de regras básicas coincide com o deles. Em outras palavras, as regras que usamos são aquelas que aparecem como padrão nas discussões filosóficas e matemáticas sobre provas de dedução natural fora dos cursos introdutórios.



\section{Regras alternativas para a quantificação}
Uma abordagem alternativa aos quantificadores é tomar como básicas as regras para $\forall$I e $\forall$E de \cref{s:BasicFOL}, além de duas regras CQ que permitem passar de $\forall \metav{x}\, \enot \metav{A}$ para $\enot \exists \metav{x}\, \metav{A}$ e vice-versa.\footnote{Warren Goldfarb segue essa linha em \textit{Deductive Logic}, 2003, Hackett Publishing Co.}

Tomando apenas essas regras como básicas, poderíamos ter derivado as regras $\exists$I e $\exists$E apresentadas em \cref{s:BasicFOL}. Derivar a regra $\exists$I é bastante simples. Suponha que $\metav{A}$ contenha o nome $\metav{c}$, não contenha nenhuma ocorrência da variável $\metav{x}$ e queiramos fazer o seguinte:
\begin{fitchproof}
	\have[m]{a}{\metav{A}(\ldots \metav{c} \ldots \metav{c}\ldots)}
	\have[k]{c}{\exists \metav{x}\, \metav{A}(\ldots \metav{x} \ldots \metav{c}\ldots)}
\end{fitchproof}
Isso ainda não é permitido, pois, nesse novo sistema, não temos a regra $\exists$I. Podemos, no entanto, oferecer o seguinte:
\begin{fitchproof}
	\have[m]{a}{\metav{A}(\ldots \metav{c} \ldots \metav{c}\ldots)}
	\open
		\hypo{nEna}{\enot \exists \metav{x}\, \metav{A}(\ldots \metav{x} \ldots \metav{c}\ldots)}\AS
		\have{Ana}{\forall \metav{x}\, \enot \metav{A}(\ldots \metav{x} \ldots \metav{c}\ldots)}\cq{nEna}
		\have{nAc}{\enot\metav{A}(\ldots \metav{c} \ldots \metav{c}\ldots)}\Ae{Ana}
		\have{red}{\ered}\ne{nAc, a}
	\close
	\have{nnEa}{\exists \metav{x}\, \metav{A}(\ldots \metav{x} \ldots \metav{c}\ldots)}\ip{nEna-red}
\end{fitchproof}\noindent
Derivar a regra $\exists$E é consideravelmente mais sutil. Isso ocorre porque a regra $\exists$E tem uma restrição importante (assim como a regra $\forall$I), e precisamos garantir que a estamos respeitando. Suponha, então, que estejamos em uma situação na qual \emph{queremos} fazer o seguinte:
\begin{fitchproof}
	\have[m]{ExA}{\exists \metav{x}\, \metav{A}(\ldots \metav{x} \ldots \metav{x}\ldots)}
	\open
		\hypo[i]{Ac}{\metav{A}(\ldots \metav{c} \ldots \metav{c}\ldots)}\AS
		\have[j]{B}{\metav{B}}
	\close
	\have[k]{end}{\metav{B}}
\end{fitchproof}\noindent
onde $\metav{c}$ não ocorre em nenhuma suposição não descarregada, nem em $\metav{B}$, nem em $\exists \metav{x}\, \metav{A}(\ldots \metav{x} \ldots \metav{x}\ldots)$. Normalmente, teríamos permissão para usar a regra $\exists$E; mas aqui não estamos supondo que temos acesso a essa regra como regra básica. Ainda assim, poderíamos oferecer a seguinte derivação mais complicada:

\begin{fitchproof}
	\have[m]{ExA}{\exists \metav{x}\, \metav{A}(\ldots \metav{x} \ldots \metav{x}\ldots)}
	\open
		\hypo[i]{Ac}{\metav{A}(\ldots \metav{c} \ldots \metav{c}\ldots)}\AS
		\have[j]{B}{\metav{B}}
	\close
	\have[k]{condi}{\metav{A}(\ldots \metav{c} \ldots \metav{c}\ldots) \eif \metav{B}}\ci{Ac-B}
	\open
		\hypo{nB}{\enot \metav{B}}\AS
		\have{nAc}{\enot \metav{A}(\ldots \metav{c} \ldots \metav{c}\ldots)}\mt{condi, nB}
		\have{AxnA}{\forall \metav{x} \enot \metav{A}(\ldots \metav{x} \ldots \metav{x}\ldots)}\Ai{nAc}
		\have{nEA}{\enot \exists \metav{x} \metav{A}(\ldots \metav{x} \ldots \metav{x}\ldots)}\cq{AxnA}
		\have{red2}{\ered}\ne{nEA, ExA}
	\close
	\have{nnB}{\metav{B}}\ip{nB-red2}
\end{fitchproof}\noindent
Podemos usar $\forall$I na linha $k+3$ porque $\metav{c}$ não ocorre em nenhuma suposição não descarregada nem em $\metav{B}$. As entradas das linhas $k+4$ e $k+1$ se contradizem, porque $\metav{c}$ não ocorre em $\exists \metav{x}\, \metav{A}(\ldots \metav{x} \ldots \metav{x} \ldots)$.

Com essas regras derivadas em mãos, poderíamos prosseguir e derivar as duas regras CQ restantes, exatamente como em \cref{s:DerivedFOL}.

Então, por que começamos com todas as regras dos quantificadores como básicas e depois derivamos as regras CQ?

Nossa primeira razão é que parece mais intuitivo tratar os quantificadores em pé de igualdade, dando-lhes suas próprias regras básicas de introdução e eliminação.

Nossa segunda razão está relacionada à discussão das regras alternativas para a negação. Nas derivações das regras $\exists$I e $\exists$E que oferecemos nesta seção, invocamos IP. Mas, como mencionamos anteriormente, IP é uma regra controversa. Portanto, se quisermos passar para um sistema que abandone IP, mas que ainda permita usar quantificadores existenciais, desejaremos tomar como básicas as regras de introdução e eliminação dos quantificadores e tomar as regras CQ como derivadas. (De fato, em um sistema sem IP, LEM e DNE, seremos \emph{incapazes} de derivar a regra CQ que passa de $\enot \forall \metav{x}\, \metav{A}$ para $\exists \metav{x}\, \enot \metav{A}$.)
